\subsection{Smoothing correction, inversion, and exact endpoints}
With the full orthogonal eigenbasis of the penalty,
$Q_d=V\operatorname{diag}(\delta_1,\ldots,\delta_n)V^\top$,
the $n-d$ positive eigenvalues are $\delta_j=\sigma_j(D_d)^2$,
and the remaining $d$ are zero. Consequently, the spectral
eigenvalues of $H_\lambda=(I+\lambda Q_d)^{-1}$ are
\begin{equation}\label{eq:h-alpha}
\alpha_j(\lambda)=
\begin{cases}(1+\lambda\delta_j)^{-1},&1\le j\le n-d,\\
1,&n-d<j\le n.\end{cases}
\end{equation}
This exposes both the exact limiting polynomial directions and
the scalar attenuation of the penalized directions. If
$\nu(\lambda)=\operatorname{tr}(H_\lambda)
=d+\sum_{j=1}^{n-d}\alpha_j(\lambda)$, the original
trace-based Guerrero index
\begin{equation}\label{eq:original-scale}
S_G(\lambda)=1-\frac{\nu(\lambda)}n
\quad\text{ranges from }0\text{ to }1-\frac dn.
\end{equation}
The upper bound is $1-d/n$, \emph{not} $1-n/d$.
The library also uses a normalized correction:
\begin{equation}\label{eq:normalized-scale}
S(\lambda)=\frac{S_G(\lambda)}{1-d/n}
=\frac{n-\nu(\lambda)}{n-d}\in[0,1].
\end{equation}
The map is strictly increasing for $\lambda\in[0,\infty)$
because
\begin{equation}\label{eq:s-derivative}
S'(\lambda)=\frac{1}{n-d}\sum_{j=1}^{n-d}
\frac{\delta_j}{(1+\lambda\delta_j)^2}>0.
\end{equation}
The exact endpoints are $S(0)=0$ and $S(\infty)=1$;
for any interior target, a one-dimensional bracketed solver
inverts this map. Equal normalized smoothness makes trace scales
comparable across $d$; it does not make the underlying trend
geometries or roughness penalties identical.

\subsection{Where the computational acceleration comes from}
For fixed $\mu$ put $p_t=\mu\binom td$ and
$c=V^\top(y-p)$. One common projection computes $c$; after
that, for each candidate $\lambda$,
\begin{align}\label{eq:fast-profiles}
\operatorname{tr}(H_\lambda)&=d+\sum_{j=1}^{n-d}\alpha_j(\lambda),\\
\|y-\widehat{\tau}_{\lambda,\mu}\|^2
&=\sum_{j=1}^{n-d}(1-\alpha_j(\lambda))^2c_j^2.
\end{align}
Thus a sweep of $M$ values can obtain trace, normalized smoothness,
GCV, AICc, BIC and related scalar profiles in $O(Mn)$
\emph{after} the shared eigendecomposition and $O(n^2)$
projection. A full reconstructed trend remains $O(n^2)$ per
penalty, and exact ordinary CV leverage diagnostics also involve
the spectral basis. The full dense SVD has $O(n^3)$
preprocessing time and $O(n^2)$ storage. This decomposition/reuse
trade-off explains the potential responsiveness of repeated-criteria
explorations without claiming unconditional superiority to
banded linear solvers \cite{weinert2007wh,biessy2026wh}.

Large batches of spectral attenuation operations are naturally
vectorizable. Dense linear algebra can also benefit from optimized
multithreaded BLAS on CPUs. In principle, a GPU backend could
parallelize matrix decompositions and batched matrix operations,
but host/device transfer, synchronization, and initialization
costs may negate a benefit for short series. We found
\texttt{numpy} and \texttt{scipy.linalg.svd} in the package
currently reviewed, not a released GPU solver; therefore no GPU
speedup, power saving, or measured device execution is reported.
A future benchmark should record CPU and GPU models, precision,
BLAS and accelerator versions, transfer time, warmups, memory, and
the number of candidate penalties.
